\documentclass[12pt,a4paper]{article}

\usepackage[T1]{fontenc}
\usepackage[utf8]{inputenc}
\usepackage[brazil]{babel}
\usepackage{graphicx}
\usepackage{geometry}
\usepackage{listings}
\usepackage{listingsutf8}
\usepackage{xcolor}
\usepackage{setspace}
\usepackage{float}
\usepackage{tikz}
\usepackage{tabularx}

\usetikzlibrary{positioning, shapes.geometric}
\onehalfspacing
\geometry{margin=2.5cm}

\lstset{
    language=SQL,
    basicstyle=\ttfamily\small,
    keywordstyle=\color{blue},
    commentstyle=\color{gray},
    stringstyle=\color{red},
    breaklines=true,
    showstringspaces=false
}

\newcommand{\montarinicio}[3]{
    % =========================
    % CAPA
    % =========================
    \begin{titlepage}
        \centering

        {\large Faculdade Anhanguera}\\[0.5cm]
        {\large Tecnólogo em DevOps}\\[3cm]

        {\Large \textbf{#1}}\\[2cm]

        {\large Gabriel Blank Stift Mousquer}\\[1cm]

        {\large Ivoti/RS}\\
        {\large 2026}

    \end{titlepage}

    % =========================
    % FOLHA DE ROSTO
    % =========================
    \begin{titlepage}
        \centering

        {\large Faculdade Anhanguera}\\[0.5cm]
        {\large Tecnólogo em DevOps}\\[2cm]

        {\Large \textbf{#1}}\\[1.5cm]

        \begin{flushright}
            Trabalho apresentado à disciplina de \\
            \textbf{#2} \\
            Professor: #3 \\
        \end{flushright}

        \vfill

        {\large Gabriel Blank Stift Mousquer}\\
        Matrícula: 2024201126\\

        \vfill

        {\large Ivoti/RS}\\
        {\large 2026}

    \end{titlepage}

    % =========================
    % SEÇÃO: TOC
    % =========================
    \tableofcontents
    \newpage
}


\begin{document}

\montarinicio
    {PROPOSIÇÕES LÓGICAS}
    {Lógica e Matemática Computacional}
    {Adriane Aparecida Loper}

%%%%%%%%%%%%%%%%%%%%%%%%%%%%%%%%%%%%%%%%%%%%%%%%%%%%%%%%%%%%%%%%%%%%%%%%%%%%%%%%
% SEÇÃO: TABELAS VERDADE
%%%%%%%%%%%%%%%%%%%%%%%%%%%%%%%%%%%%%%%%%%%%%%%%%%%%%%%%%%%%%%%%%%%%%%%%%%%%%%%%

\section{Tabelas Verdade}

Nesta seção, são apresentadas as tabelas-verdade das operações lógicas fundamentais: conjunção, disjunção e negação. Para cada uma delas, são exibidos também seus respectivos circuitos lógicos e uma implementação na sintaxe do Python. A combinação desses operadores permite a construção de portas lógicas como \textbf{NAND} e \textbf{NOR}, consideradas fundamentais por serem funcionalmente completas, isto é, capazes de representar qualquer função lógica e, portanto, implementar qualquer circuito digital.

\subsection{Conjunção}

A conjunção ($P \land Q$) é uma operação lógica que resulta em verdadeiro apenas quando ambas as proposições são verdadeiras simultaneamente. Em eletrônica digital, é implementada pela porta AND. Além de sua interpretação lógica, é amplamente utilizada na programação em operações bit a bit (bitwise), especialmente em máscaras de bits, onde permite selecionar ou verificar estados específicos em representações binárias.

\begin{figure}[h]
\centering
\begin{minipage}{0.30\textwidth}
    \centering
    \[
    \begin{array}{c|c|c}
        P & Q & P \land Q \\
        \hline
        V & V & V \\
        V & F & F \\
        F & V & F \\
        F & F & F \\
    \end{array}
    \]
\end{minipage}
\begin{minipage}{0.30\textwidth}
    \centering
    \begin{circuitikz}
        \draw
        (2,0.5) node[and port] (and1) {}

        (0,1) node[left]{P} -- (and1.in 1)
        (0,0) node[left]{Q} -- (and1.in 2)

        (and1.out) -- (2.5,0.5) node[right]{Saída};
    \end{circuitikz}
\end{minipage}
\begin{minipage}{0.35\textwidth}
    \centering
    \begin{lstlisting}[language=Python, breaklines=true, columns=fullflexible, basicstyle=\ttfamily\small]
    def conjuncao(P, Q):
        return P and Q
    \end{lstlisting}
\end{minipage}
\caption{Operação de conjunção em tabela-verdade, circuito lógico e linguagem Python}
\end{figure}

\subsection{Disjunção}

A disjunção ($P \lor Q$) é uma operação lógica que resulta em verdadeiro quando pelo menos uma das proposições é verdadeira. Em eletrônica digital, é implementada pela porta OR. Na programação, também aparece em operações bit a bit (bitwise), sendo utilizada para combinar ou ativar bits específicos em representações binárias, como em máscaras de controle e manipulação de flags.

\begin{figure}[h]
\centering
\begin{minipage}{0.30\textwidth}
    \centering
    \[
    \begin{array}{c|c|c}
        R & S & R \lor S \\
        \hline
        V & V & V \\
        V & F & V \\
        F & V & V \\
        F & F & F \\
    \end{array}
    \]
\end{minipage}
\begin{minipage}{0.30\textwidth}
    \centering
    \begin{circuitikz}
        \draw
        (2,0.5) node[or port] (or1) {}

        (0,1) node[left]{P} -- (or1.in 1)
        (0,0) node[left]{Q} -- (or1.in 2)

        (or1.out) -- (2.5,0.5) node[right]{Saída};
    \end{circuitikz}
\end{minipage}
\begin{minipage}{0.35\textwidth}
    \centering
    \begin{lstlisting}[language=Python, breaklines=true, columns=fullflexible, basicstyle=\ttfamily\small]
    def disjuncao(R, S):
        return R or S
    \end{lstlisting}
\end{minipage}
\caption{Operação de disjunção em tabela-verdade, circuito lógico e linguagem Python}
\end{figure}

\subsection{Negação}

A negação ($\neg P$) é uma operação lógica que inverte o valor de verdade de uma proposição, transformando verdadeiro em falso e vice-versa. Em eletrônica digital, é implementada pela porta NOT, sendo fundamental na construção de circuitos lógicos mais complexos e na aplicação de transformações algébricas. Na programação, também está presente em operações bit a bit (bitwise), onde inverte os bits de uma representação binária.

\begin{figure}[h]
\centering
\begin{minipage}{0.30\textwidth}
    \[
    \begin{array}{c|c}
        T & \neg T \\
        \hline
        V & F \\
        F & V \\
    \end{array}
    \]
\end{minipage}
\begin{minipage}{0.30\textwidth}
    \centering
    \begin{circuitikz}
        \draw
        (0,0) node[left]{T} -- (1,0)
        (1.5,0) node[not port] (not1) {}
        (not1.in) -- (1,0)
        (not1.out) -- (2.5,0) node[right]{Saída};
    \end{circuitikz}
\end{minipage}
\begin{minipage}{0.35\textwidth}
    \centering
    \begin{lstlisting}[language=Python, breaklines=true, columns=fullflexible, basicstyle=\ttfamily\small]
    def negacao(T):
        return not T
    \end{lstlisting}
\end{minipage}
\caption{Operação de negação em tabela-verdade, circuito lógico e linguagem Python}
\end{figure}

%%%%%%%%%%%%%%%%%%%%%%%%%%%%%%%%%%%%%%%%%%%%%%%%%%%%%%%%%%%%%%%%%%%%%%%%%%%%%%%%
% SEÇÃO: LEIS DE DE MORGAN
%%%%%%%%%%%%%%%%%%%%%%%%%%%%%%%%%%%%%%%%%%%%%%%%%%%%%%%%%%%%%%%%%%%%%%%%%%%%%%%%

\section{Leis de De Morgan}

Nesta seção, são apresentadas as Leis de De Morgan, que descrevem como a negação de proposições compostas afeta os operadores lógicos. Essas leis estabelecem que a negação de uma conjunção equivale à disjunção das negações, e vice-versa. Além de seu papel teórico na lógica proposicional, elas são fundamentais na eletrônica digital, pois permitem a transformação e simplificação de circuitos, especialmente no uso de portas \textbf{NAND} e \textbf{NOR}.

\begin{table}[h]
\centering
\[
\begin{array}{c|c|c|c}
U & V & \neg(U \land V) & \neg U \lor \neg V \\
\hline
V & V & F & F \\
V & F & V & V \\
F & V & V & V \\
F & F & V & V \\
\end{array}
\]
\caption{Tabela verdade das leis de De Morgan}
\end{table}

\begin{figure}[h]
\centering
\begin{circuitikz}

% =========================
% ESQUERDA: ¬(U ∧ V)
% =========================

% entradas
\node (U1) at (0,2) {U};
\node (V1) at (0,0) {V};

% AND
\node[and port] (AND1) at (2,1) {};

% conexões AND
\draw (U1) -- (AND1.in 1);
\draw (V1) -- (AND1.in 2);

% NOT
\node[not port] (NOT1) at (3,1) {};

\draw (AND1.out) -- (NOT1.in);

% saída
\draw (NOT1.out) -- (4,1) node[right] {Saída};

% =========================
% SEPARAÇÃO
% =========================

\draw[dashed] (6,3) -- (6,-2);

% =========================
% DIREITA: ¬U ∨ ¬V
% =========================

% entradas
\node (U2) at (7,2) {U};
\node (V2) at (7,0) {V};

% NOTs
\node[not port] (NOTU) at (8,2) {};
\node[not port] (NOTV) at (8,0) {};

\draw (U2) -- (NOTU.in);
\draw (V2) -- (NOTV.in);

% OR
\node[or port] (OR1) at (10.5,1) {};

\draw (NOTU.out) -- (OR1.in 1);
\draw (NOTV.out) -- (OR1.in 2);

% saída
\draw (OR1.out) -- (11,1) node[right] {Saída};

\end{circuitikz}
\caption{Circuitos lógicos equivalentes pelas leis de De Morgan}
\end{figure}

\newpage

%%%%%%%%%%%%%%%%%%%%%%%%%%%%%%%%%%%%%%%%%%%%%%%%%%%%%%%%%%%%%%%%%%%%%%%%%%%%%%%%
% SEÇÃO: TAUTOLOGIA, CONTRADIÇÃO E CONTINGÊNCIA
%%%%%%%%%%%%%%%%%%%%%%%%%%%%%%%%%%%%%%%%%%%%%%%%%%%%%%%%%%%%%%%%%%%%%%%%%%%%%%%%

\section{Tautologia, Contradição e Contingência}

Nesta seção, são analisadas proposições compostas com o objetivo de classificá-las como tautologia, contradição ou contingência. Para isso, são construídas tabelas-verdade que permitem verificar o comportamento lógico das expressões W, X e Y para todas as combinações possíveis de valores.

\subsection{Tautologia}

A proposição W ($P \lor \neg P$) é sempre verdadeira, independentemente do valor de $P$, sendo classificada como \textbf{tautologia}.

\begin{table}[h]
\centering
\[
\begin{array}{c|c|c|c|c|c}
P & \neg P & P \lor \neg P \\
\hline
V & F & V \\
F & V & V \\
\end{array}
\]
\caption{Tabela verdade da expressão $P \lor \neg P$}
\end{table}

\subsection{Contradição}

A proposição X ($(P \land \neg P)$) é sempre falsa, pois não é possível que $P$ e sua negação sejam verdadeiros simultaneamente, caracterizando uma \textbf{contradição}.

\begin{table}[h]
\centering
\[
\begin{array}{c|c|c|c|c|c}
P & \neg P & P \land \neg P \\
\hline
V & F & F \\
F & V & F \\
\end{array}
\]
\caption{Tabela verdade da expressão $P \land \neg P$}
\end{table}

\subsection{Contingência}

A proposição Y ($(P \lor Q) \land (\neg Q \lor R)$) apresenta valores verdadeiros e falsos dependendo das combinações de $P$, $Q$ e $R$, sendo portanto uma \textbf{contingência}. Essa classificação evidencia que seu valor lógico não é fixo, dependendo diretamente das variáveis envolvidas.

\begin{table}[h]
\centering
\[
\begin{array}{c|c|c|c|c|c|c}
P & Q & R & P \lor Q & \neg Q & \neg Q \lor R & (P \lor Q) \land (\neg Q \lor R) \\
\hline
V & V & V & V & F & V & V \\
V & V & F & V & F & F & F \\
V & F & V & V & V & V & V \\
V & F & F & V & V & V & V \\
F & V & V & V & F & V & V \\
F & V & F & V & F & F & F \\
F & F & V & F & V & V & F \\
F & F & F & F & V & V & F \\
\end{array}
\]
\caption{Tabela verdade da expressão $(P \lor Q) \land (\neg Q \lor R)$}
\end{table}

\end{document}
